\cotefile{160}
\scanpage{2}
% L01
\diagram{\begin{tikzcd}[ampersand replacement=\&] X' \arrow[d,"p"'] \\ X \end{tikzcd}}
% L01–08
\uline{Lemme} Soit $f:X'\longrightarrow X$ morphisme \inserted{\unsure{loc.}} de t.f. \deleted{\unsure{un schéma|propre}} de schémas noethériens, $V'$ l'ouvert des pts de $X'$ où $f$ est étale, $Z$ une partie fermée de $X$, $\widehat X$ le complété formel de $X$ le long de $Z$, $\widehat g=(g_n)$ une ``section \inserted{formelle}'' de $X'$ sur $X$ (i.e. \deleted{section de} \inserted{syst.} coh. de sections des $X'\mathbin{\times}_X X_n$ sur les $X_n$), $Z'=g_0(\unsure{X_0|Z_0|\widehat X_\bullet})=g_0(Z)$ (qui est un ss-schéma \deleted{\unsure{fermé}} de $X'$ \deleted{\illegible{1}}, avec $p|Z': Z'\simeq Z$).

% L09
\deleted{\unsure{\textcircled{2}|\textcircled{1}}} On pose $U_Z=g_0^{-1}(\deleted{\unsure{Z}}V')\xleftarrow{\sim}V'\cap Z'$, \textcircled{1} \deleted{\unsure{a}} On \deleted{\unsure{suppose}} \inserted{considère} \deleted{On suppose} les conditions a) $+$ b)

% L10
\margin{\uline{\unsure{que}}}
% L10–12
\[
\left\{\begin{aligned}
&\text{a) \unsure{les|La} \inserted{famille des}}\quad X_n\cap U_Z=g_n^{-1}(V')\\
&\quad\inserted{{}\xleftarrow{\sim}\text{$n^{\text{ième}}$ vois. inf. $V'_n$ de $V'$ \deleted{\unsure{en}} le long de $V'\cap Z'$}}\\
&\qquad\text{est sch. dense dans }\widehat X\\
&\text{b) }X'\text{ est l'adh. sch. des la famille des }g_n(X_n)\subset X'_n
\end{aligned}\right.
\]
% L13
Cette \deleted{\unsure{question}} \inserted{conjonction} implique \deleted{\unsure{que}}

% L14
c) La famille des $V'_n$ \inserted{(ou encore $\widehat V'$)} est sch. dense dans $X'$

% L15–17
\textcircled{2} Ceci posé, \deleted{$p$ est} \inserted{sous la \unsure{seule} condition c),} on a $V'\supset Z$, i.e. $p$ est étale aux pts de $Z$ (et il existe \inserted{donc} une partie ouverte $U'\mathrel{\unsure{\in|\subset}}V'$ contenant $Z$, tel que $U'\mathbin{\times}_X Z\unsure{\Rightarrow|=\rightarrow|\xrightarrow{\sim}}Z'\xrightarrow{\sim}Z$).

% L18–20
\uline{Dém} \deleted{OPS $X$ affine, car l'hyp. \inserted{est} stable par localisation ouverte sur $X$} La condition c) équivaut à la conjonction des conditions

% L21–22
\[
\begin{aligned}
c_1)&\ V'\text{ sch. dense dans }X'\\
c_2)&\ \widehat V'\text{ sch. dense dans }V',\quad\text{i.e.}
\end{aligned}
\]
% archive sheet number: 1
\scanpage{3}
% L01–02
$\mathbb F_p$, $G_a=G_{a\,\mathbb F_p}$, $n\in\mathbb N$, $X_n$ schéma des sous-groupes finis de rang $p^n$ de $G_a$,
% L02–04
\[
\begin{aligned}
\mathbb E^n=E^n_{\mathbb F_p}&\xrightarrow{\varphi_n}X_n\\
a=(a_0,\ldots,a_{n-1})&\longmapsto\operatorname{Ker}(f_a)\\
f_a:G_a\longrightarrow G_a &\qquad f_a(x)=x^{p^n}+a_{n-1}x^{p^{n-1}}\ \text{---}+a_1x^p+a_0x
\end{aligned}
\]

% L05
\uline{Lemme} $\varphi_n$ est une \deleted{monomorphisme} \inserted{imm. fermée} surjecti\deleted{\unsure{f}}ve, donc induit $\mathbb E^n\xrightarrow{\sim}X_{n\,\unsure{\mathrm{réd}|\mathrm{red}}}$

% L06–07
a) surjectif : il suffit de prouver $k^n\longrightarrow X_n(k)$ surj. si $k$ Alg clos. Or si $\deleted{\unsure{K}}H\subset G_{a\,k}$, $H\in X_n(k)$ i.e. $H$ fini de rang $p^n$, on a
% L08–09
\[
H\simeq\underset{H^0}{\underbrace{\alpha_{p^r}}}\mathbin{\underset{\cdot}{\times}}
\underset{\uparrow\,H\deleted{\unsure{(}}(k)_k}{\left(\mathbb Z/p\right)^s_k}
\qquad r+s=n .
\]
% L08–09
Soit $\xi_1\ldots\xi_s$ base de $H(k)$ sur $\mathbb F_p$, On a $H\deleted{\unsure{Z}}(k)_k=\operatorname{Ker}(\deleted{\unsure{g}}f_a)$, $a\in\mathbb E^s_k$,
% L10
\margin{justifier}
% L10–11
\[
\left\{\begin{aligned}
&\alpha=(\alpha_0,\ldots,\alpha_{s-1})\\
&f_{\unsure{a|\alpha}}(\unsure{x})=x^{p^s}+\alpha_{s-1}x^{p^{s-1}}\ \text{---}+\alpha_1x^p+\alpha_0x
\end{aligned}\right.
\]

% L12
Considérons alors
% L12–14
\[
\begin{aligned}
&f:G_a\mathbin{\text{---}}G_a\qquad f=f_a\operatorname{Frob}^{r}\\
&f(x)=x^{p^{r+s}}+\underset{\substack{\shortparallel\\a_{r+s-1}}}{\alpha_{s-1}}x^{p^{r+s-1}}\ \text{---}
+\underset{\substack{\shortparallel\\a_r}}{\alpha_0}x^{p^r}\\
&\hspace{15em}a_{r-1}=\ldots=a_0=0
\end{aligned}
\]
% L15
\[
\text{On a}\qquad\operatorname{Ker} f=f_a^{-1}\bigl(\underset{\substack{\shortparallel\\\operatorname{Ker}\operatorname{Frob}^r}}{\alpha_{p^r}}\bigr)\supset H(k)_k\mathbin{\cdot}\alpha_{p^r}
\]
% L16–17
et comme les 2 mbs sont de $=$ rang $p^{r+s}=p^n$, l'incl. $\supset$ est égalité.

% L18–20
b) \uline{Mono} Soient \deleted{$a,b\in\unsure{E}$} $k$ $\mathbb F_p$-Alg., $a,b\in\mathbb E^n$ tels que $\varphi_n(a)=\varphi_n(b)$ i.e. $\operatorname{Ker}f_a=\operatorname{Ker}f_b$, prouvons $f_a=f_b$ (i.e. $a=b$) \deleted{\unsure{1\textsuperscript{er} Cas} \unsure{$k$ corps alg. clos}}
% L21–22
\diagram{\begin{tikzcd}[ampersand replacement=\&,column sep=small]
H \arrow[r,phantom,"\subset" description] \& G_{a\,k}\arrow[rr,shift left=2pt,"f_a"]\arrow[rr,shift right=2pt,"f_b"']\arrow[dr] \& \& G_{a\,k}\\
\& \& G_{a\,k}/H\arrow[ur,shift left=2pt,"\unsure{\sim}"]\arrow[ur,shift right=2pt,"\sim"'] \& g_a,g_b
\end{tikzcd}}
% L23–24
$g_b=ug_a$ donc $f_b=uf_a$, $u\in\operatorname{Aut}(G_{a\,k})$. Il faut prouver $u=\mathrm{id}_{G_a}$. Clair si $k$ réduit car $\operatorname{Aut}(G_a)\simeq k^*$.
% archive sheet number: 2
\scanpage{4}
% L01–02
\deleted{Cas général \uline{OPS} $k$ noeth., et que $J$ idéal de carré nul, $k_0=k/J$, \unsure{$a_0=b_0$} $u_0={}$}

% L03–04
\uline{Lemme} Soit $a\in\deleted{\unsure{E}}k^n$, $u\in\operatorname{Aut}(G_{a\,k})$ tel que $u\circ f_a$ soit de la forme $f_b$ ($b\in k^n$), alors $u=\mathrm{id}_{G_{a\,k}}$.

% L05–06
\deleted{Cas général OPS $k$ noeth., $J$ idéal de carré nul et $u$ tel que $u_0=\mathrm{id}_{G_{a\,k_0}}$ ($k_0=k/J$). On a}
% L07
\[
u(x)=c_0x+c_1x^p+c_2x^{p^2}\ \text{---}+c_Nx^{p^N},\quad \unsure{c_i|a_i}\in J
\]
% L08–12
\[
\begin{aligned}
u\circ f_a(x)&\mathrel{{:}}c_0(x^{p^n}+a_{n-1}x^{p^{n-1}}\unsure{+|-}a_{n-2}x^{p^{n-2}}\ \text{---}+a_1x^p+a_0x)\\
&+c_1(x^{p^{n+1}}+a_{n-1}^{p}x^{p^n}\ \text{--}\qquad\text{--}+a_0^p x^p)\\
&+c_2(x^{p^{n+2}}+a_{n-1}^{p^2}x^{p^{n+1}}+\qquad +a_0^{p^2}x^{p^2})\\
&\text{---}\\
&+c_N({\unsure{x|\alpha}}^{p^{n+N}}+a_{n-1}^{p^N}x^{p^{n+N-1}}+\text{---}\qquad \unsure{-|+}a_0^{p^N}x^{p^N})
\end{aligned}
\]
% L13–15
Le coeff\textsuperscript{t} de $x^{p^{n+N}}$ dans l'expression totale est $c_N x^{p^{n+N}}$, donc \deleted{\unsure{pour}} si $N\geq 1$, on doit avoir $c_N=0$. Donc de proche en proche $c_{N-1}=0\ldots c_1=0$.

% L16–17
Donc $u(x)=c_0(x)=c_0 x^{p^n}\ \text{---}$, et il faut avoir $c_0=1$ cqfd.

% L18–19
c) $\varphi_n$ \deleted{\unsure{est}} \uline{propre} Critère valuatif, il faut prouver \deleted{\uline{\unsure{Corollaire|Conclusion}} \unsure{$X_{n\,\mathrm{réd}}$ est intègre}}

% L20–24
\uline{Lemme} Soit $S$ \inserted{$=\operatorname{Spec}(V)$} trait sur $\mathbb F_p$, \inserted{pt gén. $y$} $H\subset G_{a\,S}$, \inserted{$H\in X_n(S)$ i.e. $H$ plt fini $/S$ de rg $p^n$,} supposons que $H_y=\operatorname{Ker} f_{a_y}$, $\unsure{a_y|a|a_\eta}\in K^n$ ($K=k(y)$), alors $a_y\in V^n$ et \deleted{\unsure{$H=$} \unsure{$\operatorname{Ker}(f_a:G_{a\,S}$}} $H=\operatorname{Ker}(f_a:G_{a\,S}\mathbin{\text{---}}G_{a\,S})$.
% archive sheet number: 3
\scanpage{5}
% L01–05
Supposons que $a_y=(a_0,a_1,\text{--},a_{n-1})$ soit $\notin V^n$, \deleted{\unsure{donc que}} soit $\unsure{\nu|v}=-\operatorname{Inf}\operatorname{val}(a_i)$, $\pi$ uniformisante de $V$, Soit $b=\pi^\nu a=(\pi^\nu a_0,\text{--},\pi^\nu a_{n-1})$, de sorte que que les $b_i=\pi^\nu a_i$ \deleted{soit} ($0\leq i\leq n-1$) sont $\in V$ et l'un est une unité. On a, si
% L06–07
\[
\begin{aligned}
&f(x)=\pi^\nu x^{p^n}+b_{n-1}x^{p^{n-1}}\ \text{---}+b_1x^p+b_0x\\
&f:G_{a\,S}\mathbin{\text{---}}G_{a\,S}
\end{aligned}
\]
% L08
\[ H_y=\operatorname{Ker} f_y \]
% L09–15
Par ailleurs $f$ est \deleted{\unsure{épi}} fid plat -- car on le voit fibre {:} fibre, et sur la fibre spéciale \unsure{ceci|cela} provient du fait que $f_s$ n'est pas nul (car un des coeff. $b_{is}$ est $\ne 0$). Donc $\operatorname{Ker} f$ est un sous-schéma en groupes \inserted{fermé} plat de $G_{a\,S}$, de fibre gén. $H_y$, donc coïncide avec $H$ \deleted{Donc} (qui \inserted{comme lui} est l'adh. sch. de $H_y$ dans $G_{a\,S}$).

% L16–19
Donc $\operatorname{Ker} f_s$ est de rang $p^n$, or il est de rang $p^{n-i}$ si $b_{n-i,s}$ \deleted{\unsure{qui}} est le premier de $b_{\unsure{n-1\,s|n-1,s}}, b_{n-\deleted{\unsure{1}}2,s}$ etc qui soit non nul. C'est absurde.

% L20–23
\uline{Lemme} Soit $X_n^{\mathrm{ét}}\subset X_n$ l'ouvert correspondant aux sous-schémas en groupes étales de $G_a$. Alors $\varphi_n$ induit iso
% L21
% slanted
\margin{\uline{justifier}}
% L23
\[ \mathbb E^*\times\mathbb E^{n-1}\ \deleted{\unsure{\text{défini par}}}\xrightarrow{\sim}X_n^{\mathrm{ét}} \]
% archive sheet number: 4
\scanpage{6}
% L01–03
(on identifie $\mathbb E^*\times\mathbb E^{n-1}$ à l'ouvert de $\mathbb E^{\unsure{n}}$ dont les \deleted{\unsure{élé}} $k$-pts sont les $(a_0,\text{--},a_{n-1})\in k^n$ tels que $a_0\in k^*$ \ldots.).

% L04–06
\uline{Question} $X_n$ est-il réduit (i.e. lisse i.e. $\simeq\mathbb E^n$) ? Ou au contraire l'ens de ses pts réduits est il réduit à $X_n^{\mathrm{ét}}\simeq\mathbb E^*\times\mathbb E^{n-1}$ ?

% L07
Cas $n=1$, $X_1$ est-il réduit \inserted{i.e. lisse} au pt $0$ ?

% L08–13
Il revient au $=$ de dire que l'espace tangent de Zariski de $X_1$ en $0$ est de dim $1$ (\unsure{à présent|à priori} c'est de dim $\geq1$). \deleted{\unsure{Donc}} Il faut étudier les pts de $X_0$ à val. dans $k=\mathbb F_p[\unsure{t|T|J}]/t^2$ se réduisant suivant $0$, i.e. \unsure{c} les sous-groupes plats \inserted{$H$} de $G_{a\,k}$ tels que $H_0=\alpha_p$.

% L14–15
\diagram{\begin{tikzcd}[ampersand replacement=\&,column sep=small,row sep=small]
 H\overset{?}{\subset}\underset{\substack{\shortparallel\\G_a}}{G} \& H_0\subset G_0\arrow[r] \& \deleted{\unsure{G_0/H}}K_0\simeq G_0/H_0\\
 \& \overline H_0\subset G \& \overline K_0\simeq G/\overline H_0\simeq\overline G_0/\overline H_0
\end{tikzcd}}
% L16–20
$K=G/H$ est une \unsure{prol.|part.} de $K_0=G_0/H_0$ défini par une extension de $K_0$ par $\unsure{\mathfrak g K_0|\mathfrak L K_0|\mathbb V t_{K_0}|\mathscr L K_0|J K_0}$ (\unsure{don|dont} les splittings correspondent aux hom de $\overline K_0$ \inserted{$=G/\overline H_0$} avec $\deleted{\unsure{\overline K_0}}K\simeq G/H\deleted{\unsure{{}\simeq\overline G_0/\overline H_0}}$ \deleted{\unsure{)}} qui prolongent l'hom identique)

% L18–23
% slanted
\margin{Dégager \unsure{lemmes} pour calcul \unsure{des ces|de ces} espaces \unsure{tangents} de Zariski et obstructions à prolonger \unsure{ss-groupe} plat d'un \unsure{groupe} comm\ldots}
% L21–24
Image inverse de cette extension par $G_0\mathbin{\text{---}}K_0$ est l'obstruction à prolonger $G_0\overset{\unsure{\xi|\varepsilon}}{\mathbin{\text{---}}}K_0$ en $G\longrightarrow\deleted{\unsure{\overline K_0}}K$, or celle-ci se prolonge de façon canonique, \deleted{\unsure{en $G$}} donc cette extension
% archive sheet number: 5
\scanpage{7}
% L01–03
est \unsure{splittée|splittie}, d'où un hom \inserted{$u$} de $H_0$ dans $t_{K_0}$ qui donne l'extension de $\unsure{K_0|K}$ par $\unsure{\mathfrak g K_0|\mathfrak L K_0|\mathbb V t_{K_0}|\mathscr L K_0|J K_0}$ plus l'hom $G\longrightarrow K$ \deleted{\unsure{de}} (donc permet de réconstituer $H$)
% L04
\[ u:H_0\longrightarrow\mathfrak g K_0 \]
% L05–09
\deleted{Ici $H_0=\alpha_p$ \inserted{$\subset G_0=G_a$}, $K_0\simeq G_{a\unsure{p}}$, $\mathfrak g K_0\simeq G_a$, $\operatorname{Hom}(\alpha_p,G_a)$ et on sait $\operatorname{Hom}(\alpha_p,\unsure{\mathrm{fibr\acute e}|\mathrm{faisceau}})\xrightarrow{\sim}\operatorname{Hom}(\mathfrak g\alpha_p,\mathfrak g\unsure{\mathrm{fibr\acute e}|\mathrm{faisceau}})$ donc $\operatorname{Hom}(K_0,\mathfrak g K_0)\mathbin{\text{---}}\operatorname{Hom}(H_0,\mathfrak g K_0)$ surj, donc \deleted{l'extension} \unsure{le prolongt} $\overline K$ de $K_0$ est trivial i.e. $\simeq$}
% L09–10
On trouve l'espace tangent de Zariski :
% L10
\[ \operatorname{Hom}_{\deleted{\unsure{k}}\unsure{k\text{-gr}|\text{gr}}}(H_0,\mathfrak g K_0) \]
% L11
\begin{center}- - -\end{center}

% L12
\uline{Théorème} $\varphi_n:\mathbb E^n_{\mathbb F_p}\longrightarrow X_n$ \uline{iso}

% L13–17
i.e. Si $S$ est un schéma de car $p$ $H\subset G_{a\,S}$, ss-schéma en \unsure{gpes|groupes} fini et \deleted{\unsure{plat}} \inserted{loc. libre} de rang $p^n$, alors $\exists1\ a=(a_0,\text{--},a_{n-1})\in\mathbb E^n(S)=\Gamma(S,\mathscr O_S)^n$ tel que
% L17
\[ H=\operatorname{Ker} f_a \]
% L18
où
% L18
\[ f_a:G_{a\,S}\mathbin{\text{---}}G_{a\,S}\qquad f_a(x)=x^{p^n}+a_{n-1}x^{p^{n-1}}\ \text{--}+a_0x . \]
% L19
De plus, $H$ est étale $/S$ sss $a_0\in\Gamma(S,\mathscr O_S)^*=\Gamma(S,\mathscr O_S^*)$
% archive sheet number: 6
\scanpage{8}
% L01–02
$\mathbb E^n$ sur $\mathbb F_p$, $\Gamma_n=\operatorname{Gl}(n,\mathbb F_p)_{\mathbb F_p}\subset\operatorname{Gl}(n)_{\mathbb F_p}$ \inserted{opère sur $\mathbb E^n$}
% L03
\[ \unsure{Q|\mathbb Q}_n=\mathbb E^n/\Gamma_n=\operatorname{Spec}(A_n)\qquad A_n=\mathbb F_p[X_1,\ldots,X_n]^{\Gamma_n} \]
% L04–05
Soit $U_n\subset\mathbb E^n$ l'ouvert formé des ``syst. lin. indép. sur $\mathbb F_p$'' i.e. l'\deleted{\unsure{ens}} \unsure{plus} grand ouvert de $\mathbb E^n$ sur lequel $\Gamma_n$ opère librement, $Q_n^0=U_n/\Gamma_n$

% L06
\uline{Lemme} \deleted{On a}
% L06
{\scriptsize
\[
\deleted{\det\begin{pmatrix}X_1&X_2&\cdots&X_n\\X_1^p&X_2^p&\cdots&X_n^p\\\ldots&&&\\X_1^{p^{n-1}}&X_2^{p^{n-1}}&\cdots&X_n^{p^{n-1}}\end{pmatrix}=
\prod_{\alpha\in\mathbb F_p^{n-1}}(X_n+\sum_1^{n-1}\alpha_iX_{\unsure{n|i}})
\prod_{\alpha\in\mathbb F_p^{\unsure{n-2|n-1}}}X_{n-2}+\sum_1^{\unsure{n-2}}\alpha_iX}
\]}
% L07
On a dans $\mathbb F_p[X_1,\text{--},X_n]$
% L08–11
{\scriptsize
\[
\begin{aligned}
\underset{\substack{\shortparallel\,\mathrm{d\acute ef}\\ \Delta_n(X_1,\text{--},X_n)}}{\det\begin{pmatrix}
X_1&\cdots&X_n\\X_1^p&\cdots&X_n^p\\\ldots&&\\X_1^{p^{n-1}}&\cdots&X_n^{p^{n-1}}\end{pmatrix}}
&=\prod_{(\alpha_1,\ldots,\alpha_{n-1})\in\mathbb F_p^{n-1}}\left(X_n+\sum_1^{n-1}\alpha_i X_i\right)
\prod_{(\alpha_1,\ldots,\alpha_{n-2})\in\mathbb F_p^{n-2}}\left(X_{n-1}+\sum_1^{n-2}\alpha_iX_i\right)\ldots\\
&\qquad\ldots\left(\prod_{\alpha_1\in\mathbb F_p^1}(X_2+\alpha_1X_1)\right)X_1\\
&=\prod_{0\leq j\leq n-1}\deleted{\unsure{X_{j+1}}}\prod_{\alpha\in k^j}\left(X_{j+1}+\sum_{i=1}^j\alpha_i X_i\right)
\end{aligned}
\]}

% L12–16
\uline{NB} À un facteur $\in\mathbb F_p^*$ près, le deuxième mb est le produit des combinaisons linéaires non \deleted{\unsure{toutes}} id. nulles des $X_i$ à coeff dans $\mathbb F_p$, \deleted{\unsure{en}} \deleted{\unsure{ne pre}} prenant ces combinaisons mo\uline{dulo} \deleted{\unsure{homothétie}} \inserted{multiplication} par un $c\in\mathbb F_p^*$. C'est un polynôme \inserted{$P$} de degré
% L17
\[ 1+p\ \text{---}+p^{n-1}=\frac{p^n-1}{p-1}=\operatorname{card}P^{n}(\mathbb F_p) \]
% L18–23
donc de degré égal au premier, $Q$. Or \deleted{\unsure{les}} l'ens. des zéros de \deleted{produits} \inserted{chaque facteur $P_\lambda$ de $P$} est contenu dans celui de $Q$, \deleted{\unsure{et comme $P$}} donc \inserted{($P_\lambda$ étant premier)} $P_\lambda$ divise $Q$, donc (les $P_\lambda$ étant premiers 2 à 2 \deleted{\unsure{dist}} non associés) $Q$ est divisible par $P=\prod_\lambda P_\lambda$, donc, étant de $=$ degré, est de la forme $cP$, $c\in\mathbb F_p^*$. Reste à déterminer $c$ \ldots
% archive sheet number: 7
\scanpage{9}
% L01–02
\uline{Corollaire} Soit $k$ corps de car. $p$, $\unsure{\xi|\beta}=(\xi_1,\text{--},\xi_n)\in k^n$. Alors $\deleted{\det}\Delta(\xi_1,\text{--},\xi_n)=0$ sss $\xi_1,\text{--},\xi_n$ sont lin. dép. sur $\mathbb F_p$.

% L03
\uline{Corollaire} \[ U_n=\mathbb E_{n\,\Delta_n}=\operatorname{Spec}k[X_1,\text{--},X_n]_{\Delta_n} \]

% L04–05
Si $\xi_1\text{---}\xi_n\in\deleted{\unsure{A}}k$ anneau de car. $p$, cherchons $\deleted{\unsure{a_{n-1}}}a=(a_0,a_1,\text{--},a_{n-1})\in k^n$ tel que l'on ait
% L06
\[ \xi_i^{p^n}+a_{n-1}\xi_i^{p^{n-1}}\text{---}+a_1\xi_i^p+a_0\xi_i=0 \tag*{(*)} \]
% L07
i.e.
% L07
\[ f_a(\xi_i)=0 \qquad 1\leq i\leq n \]

% L08
\mspage{\unsure{7}}
% L08–15
C'est un syst. de $n$ équations \deleted{\unsure{différentes|diff.}} \inserted{lin. entre $a_i$} avec seconds mbs $-\xi_i^{p^n}$, dont le déterminant est $\Delta_n(\xi_1,\text{--},\xi_n)$. Sous réserve que celui-ci est inversible (i.e. les $\xi_i$ engendrent un ss-groupe \inserted{fini} étale de $G_{a\,k}$ de rang $p^n$) il y a une solution et une seule, donnée par les formules de Cramer

% L16
\[
a_0=\frac{\deleted{\unsure{1|-1}}1}{\Delta_n(\xi)}\det
\begin{pmatrix}\xi_1^{p^n}&\cdots&\xi_n^{p^n}\\\xi_1^p&\cdots&\xi_n^p\\\ldots&&\\\xi_1^{p^{n-1}}&\cdots&\xi_n^{p^{n-1}}\end{pmatrix}
=\frac{\Delta_0(\xi)}{\Delta_n(\xi)}
\]
% L17–18
\[
\begin{aligned}
&\text{--}\\
a_i&=\frac{1}{\Delta_n(\xi)}\det
\begin{pmatrix}
\xi_1&\cdots&\xi_n\\\ldots&&\\
\xi_1^{p^{i-1}}&\cdots&\xi_n^{p^{\unsure{i+1|i-1}}}\\
\xi_1^{p^n}&\cdots&\xi_n^{p^n}\\
\xi_1^{p^{i}}&\cdots&\xi_n^{p^i}\\
\ldots&&\\
\xi_1^{p^{n-1}}&\cdots&\xi_n^{p^{n-1}}
\end{pmatrix}
=\frac{\Delta_i(\xi)}{\Delta_n(\xi)}
\end{aligned}
\]
% L18
\margin{modulo signes dont je ne suis pas sur.}
% L19–20
\[
\begin{aligned}
&\text{--}\\
a_{n}&=\frac{1}{\Delta_n(\xi)}\det
\begin{pmatrix}\xi_1&\cdots&\xi_n\\\ldots&&\\
\xi_1^{p^{n-2}}&\cdots&\xi_n^{p^{n-2}}\\
\xi_1^{p^n}&\cdots&\xi_n^{p^n}\end{pmatrix}
=\frac{\Delta_{n-1}(\xi)}{\Delta_n(\xi)}
\end{aligned}
\]
% archive sheet number: 8
\scanpage{10}
% L01–03
\uline{Lemme} Prenons $k=\mathbb F_p[X_1,\text{--},X_n]$, $\xi_i=X_i$, donc les \deleted{\unsure{différents|systèmes}} $\Delta_i\in\mathbb F_p[X_1,\text{--},X_n]$ ($0\leq i\leq n$) \deleted{\unsure{donc}}, et on \deleted{Alors} a
% L04
\[ \Delta_n X_i^{p^n}+\Delta_{n-1}X_i^{p^{n-1}}\text{---}+\Delta_1X_i^p+\Delta_0X_i=0 \]
% L05–06
($1\leq i\leq n$). Les $\Delta_i$ ($0\leq i\leq n-1$) sont divisibles par $\Delta_n$, soit
% L07
\[ \text{et}\qquad \Delta_i=A_i\Delta_n \qquad 0\leq i\leq n-1 \]
% L08
\[ \text{et On a donc}\qquad \Delta_0=\pm\Delta_n^p(-1)^n\quad\text{au signe près} \]

% L09
\uline{Cor 1} \unsure{Ces|Cor}
% L09–10
\[ X_i^{p^n}+A_{n-1}\deleted{\unsure{X}}X_i^{p^{n-1}}\text{---}+A_1X_i^p\unsure{-|+} A_0X_i=0 \qquad (1\leq i\leq n) \]
% L11–12
\deleted{\unsure{et}} et les $A_i\in k[X_1,\text{--},X_n]$ ($0\leq i\leq n-1$) sont caractérisés par cette relation. \deleted{\unsure{Enfin|En effet}} De plus $A_0=\pm\Delta_n^{p-1}(-1)^n$ \deleted{\unsure{(au signe}}.

% L11–15
% margin note written sideways
\margin{\uline{Cor. 2} les $A_i$ sont \deleted{\unsure{stables}} \inserted{invariants} par l'action de $\operatorname{Gl}(n,\mathbb F_p)=\Gamma_n$ (a fortiori \deleted{et} \uline{symétriques}); $A_i$ est \uline{homogène} de degré $p^n-p^i$.}
% L13–15
\uline{Cor. 2} Revenons {:} une \deleted{$k$ qqcle. sur} \inserted{extn $k$ de} $\mathbb F_p$. Pour que l'on ait $\Delta_n(\xi)\neq0$ i.e. $\xi_1\text{--}\xi_n$ lin. indép sur $\mathbb F_p$, il f et s que $\Delta_0(\xi)\neq0$.

% L16
Considérons
% L16
\diagram{\begin{tikzcd}[ampersand replacement=\&]
\mathbb E^n \arrow[r,"A_*"] \& \mathbb E^n
\end{tikzcd}}
% L17
défini par $A_0\text{--}A_{n-1}$. Il se factorise en
% L18
\deleted{\unsure{Prop} \unsure{Cor 3} \unsure{Le $A_*$ induit un iso}}
% L19
\[ \deleted{U_n\,(=\mathbb E^n_{\Delta_n})=A_*^{-1}(\mathbb E^*\times\mathbb E^{n-1})} \]
% L20
\deleted{et $A_*$ induit un iso}
% L21–22
\diagram{\begin{tikzcd}[ampersand replacement=\&]
\mathbb E^n \arrow[rr,no head,"A_*"] \arrow[dr,no head,"\mathrm{can}" description] \& \& \mathbb E^n\\
\& \unsure{\mathbb Q|Q}_n \arrow[ur,"\psi_n"'] \&
\end{tikzcd}}
% L23
\deleted{et induit} Ceci dit.
% archive sheet number: 9
\scanpage{11}
% L01–03
\uline{Prop} \[
\psi_n^{-1}\bigl(\underbrace{\mathbb E^*\times\mathbb E^{n-1}}_{\substack{(a_0,\text{--},a_{n-1})\\\text{tels que }a_0\text{ inv.}}}\bigr)=Q_n^0
\]
% L04
et
% L05
\diagram{\begin{tikzcd}[ampersand replacement=\&]
 Q_n^0 \arrow[r,"\sim"] \& \mathbb E^*\times\mathbb E^{n-1}
\end{tikzcd}}

% L06–09
\deleted{En effet,} [Il est immédiat que $Q_n^0\simeq$ ``ss-groupes finis étales de $G_a$ de rg $p^n$'', ce qui donne une interprétation géom de ce foncteur \unsure{ouvert|couvert|constant} $\mathbb E^*\times\mathbb E^{n-1}$, en termes de noyaux d'hom $f_a$ \ldots]

% L10–11
En effet, on a un \deleted{foncteur} hom en sens inverse
% L11–12
\diagram{\begin{tikzcd}[ampersand replacement=\&]
Q_n^0 \& E^*\times\mathbb E^{n-1}\subset\mathbb E^n \arrow[l]\\
f_a \& a \arrow[l,mapsto] \& \ldots
\end{tikzcd}}

% L13–14
\uline{Thm} $\unsure{\mathbb Q|Q}_n\longrightarrow\mathbb E^n$ est un isomorphisme, en d'autres termes
% L15
\[ \mathbb F_p[X_1,\text{--},X_n]^{\Gamma_n}\simeq\mathbb F_p[A_0,\text{--},A_{n-1}] \]

% L16–19
\unsure{On|Il|Dès} sait déjà que $Q_n\mathbin{-}\mathbb E_{\deleted{\unsure{n}}}^{n}$ est birationnel, il suffit de prouver qu'il est fini \inserted{ou même propre} -- comme $Q_n$ est intègre et $E_n$ normal, on conclut que c'est un iso.
% L19–21
\deleted{Il suffit de prouver} \inserted{Ou aussi il suffit} \deleted{que $\mathbb E^n\xrightarrow{A_*}\mathbb E^n$} de prouver que $Q_n\to\mathbb E^n$ \unsure{est|$n$} est quasi-fini (car \unsure{sera|ser.} une imm. ouverte par le
% archive sheet number: 10
\scanpage{12}
% L01–03
``Main Theorem'' de Zariski) et surjectif ; et il suffit de prouver la $=$ chose pour $\mathbb E^n\xrightarrow{A_*}\mathbb E^n$, donc prouver que sur $k$ corps alg. clos, l'application
% L04
\deleted{\unsure{a)|2)} \uline{Surjectivité}}
\[ k^{\unsure{*|n}}\xrightarrow{A_*}k^n \]
% L05
est \uline{surjective} à \uline{fibres finies}.

% L06–08
1) Surjection. Soit $a=(a_0,\text{--},a_{n-1})\in k^n$. Si $a=0$, c'est l'image de $0=(\underbrace{0,\text{--},0}_{n})\in k^n$. \deleted{\unsure{Si} $\longleftarrow$}
% L08–11
Sinon, soit $a_r$ le premier parmi $a_0\text{--}a_{\unsure{n|n-1}}$ qui soit $\neq0$. Considérons $a_r\text{--}a_n\in k^{n-r+1}$, il provient de $\xi_1,\text{--},\xi_{n-r+1}$ ; on complète par des $0$ pour avoir $\xi_1\text{---}\xi_n$, ça marche, grâce au
% L12–13
\uline{Lemme} \[
\left\{\begin{aligned}
 A_0(X_1,\text{--},X_{n-1},0)&=0\\
 A_i(\deleted{\unsure{0,}}X_1,\text{--},X_{n-1},0)&=A_{i-1}(X_1,\text{--},X_{n-1})
\end{aligned}\right.
\]
pour $1\leq i\leq\unsure{n|n-1}$
% L14–15
donc
\[
 A_\alpha[X_1,\text{--},X_{n-r},\overbrace{0\text{--},0}^{\unsure{r|n-r|s}}]
 =\begin{cases}
 0&\text{si }0\leq\alpha\unsure{\leq r-1|<r|<s}\\
 A_{\alpha-r}(X_1,\text{--},X_{n-r})&\text{si }\unsure{r|s}\leq\alpha\leq n
 \end{cases}
\]

% L16–17
2) À fibres finies. En fait, on voit \unsure{tt|lt} de suite que les fibres sont les trajectoires de $\Gamma_n$.

% L18–20
[Ainsi, indépendamment du M. Th. de Zar, on voit que $Q_n\to\mathbb E^n$ est surjectif, radiciel, birationnel.

% L21–23
\uline{NB} On peut aussi regarder les invariants de $\mathbb F_p[X_1,\text{--},X_n]$ sous $\unsure{S\Gamma_n|\Sigma\Gamma_n}=\operatorname{Sl}(n,\mathbb F_p)$. Parmi ces invariants il y a $\Delta_n$ (on a $\deleted{\unsure{\Delta|\triangle}}$ en fait
% archive sheet number: 11
